\subsection{FINAL TOPOLOGIES}

PROPOSITION 6. Let X be a set, let \((\mathbf{Y}_t)_{t \in I}\) be a family of topological spaces, and for each \(t \in I\) let \(f_t\) be a mapping of \(\mathbf{Y}_t\) into X. Let \(\mathfrak{D}\) be the set of subsets U of X such that \(\overline{f}_t^1(\mathbf{U})\) is open in \(\mathbf{Y}_t\) for each \(t \in I\) ; then \(\mathfrak{D}\) is the set of open subsets of X in a topology \(\mathfrak{C}\) on X which is the final structure on X for the family \((f_t)\) (Set Theory, Chapter IV, § 2, no. 5), and in particular \(\mathfrak{C}\) is the finest topology on X for which the mappings \(f_t\) are continuous. In other words, if g is a mapping of X into a topological space Z, then g is continuous (X carrying the topology \(\mathfrak{C}\) ) if and only if each of the mappings \(g \circ f_t\) is continuous.

It is immediately verified that \(\mathfrak{D}\) satisfies the axioms \((\mathrm{O_I})\) and \((\mathrm{O_{II}})\) . {[}Set Theory, R, § 4, formulae (34) and (46){]}. We shall prove the last assertion of the proposition, which implies the other assertions by reason of the general properties of final structures (Set Theory, Chapter IV, § 2, no. 5, criterion CST 18). It is clear that the \(f_{i}\) are continuous in the topology \(\mathfrak{C}\) , by the definition of \(\mathfrak{D}\) (no. 1, Theorem 1); hence if \(g\) is continuous, so is each mapping \(g \circ f_{i}\) (no. 1, Theorem 2). Conversely, suppose that each \(g \circ f_{i}\) is continuous, and let V be an open set in Z; by hypothesis, \(\overline{f}_{i}^{1}(\overline{g}^{1}(V))\) is open in \(Y_{i}\) for each \(i \in I\) ; hence \(\overline{g}^{1}(V) \in \mathfrak{D}\) , and the proof is complete.

COROLLARY. Under the hypotheses of Proposition 6, a subset F of X is closed in the topology \(\mathfrak{C}\) if and only if \(\overline{f}_{\iota}^{1}(\mathbf{F})\) is closed in \(Y_{\iota}\) for each \(\iota \in I\) . This follows from the definition of the open sets in the topology \(\mathfrak{C}\) by taking complements.

The general properties of final structures (Set Theory, Chapter IV, § 2, no. 5, criterion CST 19) imply the following transitivity property (the direct proof of which is also straightforward):

PROPOSITION 7. Let X be a set, \((Z_{\iota})_{\iota \in I}\) a family of topological spaces, \((J_{\lambda})_{\lambda \in L}\) a partition of I, and \((Y_{\lambda})_{\lambda \in L}\) a family of sets indexed by L. Also, for each \(\lambda \in L\) let \(h_{\lambda}\) be a mapping of \(Y_{\lambda}\) into X; for each \(\lambda \in L\) and each \(\iota \in J_{\lambda}\) let \(g_{\lambda \iota}\) be a mapping of \(Z_{\iota}\) into \(Y_{\lambda}\) , and put \(f_{\iota} = h_{\lambda} \circ g_{\lambda \iota}\) . If each \(Y_{\lambda}\) carries the finest topology for which the mappings \(g_{\lambda \iota} (\iota \in J_{\lambda})\) are continuous, then the finest topology on X for which the \(f_{\iota}\) are continuous is the same as the finest topology for which the \(h_{\lambda}\) are continuous.

\paragraph*{Examples}

\begin{enumerate}
\def\labelenumi{\arabic{enumi})}
\item
  Quotient topology. Let X be a topological space, R an equivalence relation on X, Y = X/R the quotient set of X with respect to the relation R, \(\varphi: X \to Y\) the canonical mapping. The finest topology on Y for which \(\varphi\) is continuous is called the quotient of the topology of X by the relation R; we shall study it in more detail in § 3.
\item
  Greatest lower bound of a set of topologies. Every family \((\mathfrak{C}_t)_{i \in I}\) of topologies on a set \(X\) has a greatest lower bound \(\mathfrak{C}\) in the set of all topologies on \(X\) , i.e.~\(\mathfrak{C}\) is the finest of all the topologies on \(X\) which are coarser than each of the \(\mathfrak{C}_t\) . To see this we may apply Proposition 6, taking \(Y_t\) to be the set \(X\) with the topology \(\mathfrak{C}_t\) , and \(f_t\) to be the identity mapping \(Y_t \to X\) . If \(\mathfrak{D}_t\) is the set of subsets of \(X\) which are open in the topology \(\mathfrak{C}_t\) , then the set \(\bigcap_{i \in I} \mathfrak{D}_t\) is the set of subsets of \(X\) which are open in \(\mathfrak{C}\) .
\end{enumerate}

\(\mathfrak{G}\) is also called the intersection of the topologies \(\mathfrak{G}_{t}\) .

\begin{enumerate}
\def\labelenumi{\arabic{enumi})}
\setcounter{enumi}{2}
\tightlist
\item
  Sum of topological spaces. Let \((\mathbf{X}_t)_{i \in I}\) be a family of topological spaces, \(\mathbf{X}\) the set which is the sum of the \(\mathbf{X}_t\) (Set Theory, Chapter II, § 4, no. 8, Definition 8); for each \(i \in I\) , let \(j_t\) be the canonical (injective) mapping of \(\mathbf{X}_t\) into \(\mathbf{X}\) . The finest topology \(\mathfrak{G}\) on \(\mathbf{X}\) for which the mappings \(j_t\) are continuous is called the sum of the topologies of the \(\mathbf{X}_t\) , and \(\mathbf{X}\) with this topology is said to be the sum of the topological spaces \(\mathbf{X}_t\) . Let us identify each of the \(\mathbf{X}_t\) with a subset of \(\mathbf{X}\) by means of \(j_t\) ; then a set \(\mathbf{A} \subset \mathbf{X}\) is open (resp. closed) in the topology \(\mathfrak{G}\) if and only if each of the sets \(\mathbf{A} \cap \mathbf{X}_t\) is open (resp. closed) in \(\mathbf{X}_t\) . In particular, each of the \(\mathbf{X}_t\) is both open and closed.
\end{enumerate}

The following proposition generalizes the situation of Example 3:

PROPOSITION 8. Let X be a set, \((\mathbf{X}_{\lambda})_{\lambda \in \mathbf{L}}\) a family of subsets of X. Suppose each \(X_{\lambda}\) carries a topology \(C_{\lambda}\) such that, for each pair of indices \((\lambda, \mu)\) :

\begin{enumerate}
\def\labelenumi{\arabic{enumi})}
\item
  \(\mathbf{X}_{\lambda} \cap \mathbf{X}_{\mu}\) is open (resp. closed) in each of the topologies \(\mathfrak{C}_{\lambda}, \mathfrak{C}_{\mu}\) .
\item
  The topologies induced on \(\mathbf{X}_{\lambda} \cap \mathbf{X}_{\mu}\) by \(\mathfrak{C}_{\lambda}\) and \(\mathfrak{C}_{\mu}\) coincide. Let \(\mathfrak{C}\) be the finest topology on \(\mathbf{X}\) for which the injections \(j_{\lambda}: \mathbf{X}_{\lambda} \to \mathbf{X}\) are continuous. Then, for each \(\lambda \in \mathbf{L}\) , \(\mathbf{X}_{\lambda}\) is open (resp. closed) in \(\mathbf{X}\) in the topology \(\mathfrak{C}\) , and the topology induced by \(\mathfrak{C}\) on \(\mathbf{X}_{\lambda}\) coincides with \(\mathfrak{C}_{\lambda}\) .
\end{enumerate}

In view of Proposition 6 and its corollary, it is enough to show that for each \(\lambda\) and each subset \(A_{\lambda}\) of \(X_{\lambda}\) the following statements are equivalent:

\begin{enumerate}
\def\labelenumi{(\roman{enumi})}
\item
  \(\mathbf{A}_{\lambda}\) is open (resp. closed) in the topology \(\mathfrak{C}_{\lambda}\) .
\item
  For each \(\mu \in L\) , \(A_{\lambda} \cap X_{\mu}\) is open (resp. closed) in the topology \(\mathfrak{G}_{\mu}\) .
\end{enumerate}

It is clear that (ii) implies (i) by taking \(\mu = \lambda\) . Conversely, if (i) is satisfied, then \(A_{\lambda} \cap X_{\mu}\) is open (resp. closed) in \(X_{\lambda} \cap X_{\mu}\) for the topology \(\mathcal{C}_{\lambda\mu}\) induced on \(X_{\lambda} \cap X_{\mu}\) by \(\mathcal{C}_{\lambda}\) ; but \(\mathcal{C}_{\lambda\mu}\) is also the topology induced on \(X_{\lambda} \cap X_{\mu}\) by \(\mathcal{C}_{\mu}\) ; hence \(A_{\lambda} \cap X_{\mu}\) is also the intersection of \(X_{\lambda} \cap X_{\mu}\) with a subset \(B_{\mu}\) of \(X_{\mu}\) which is open (resp. closed) in the topology \(\mathcal{C}_{\mu}\) ; since \(X_{\lambda} \cap X_{\mu}\) is open (resp. closed) in \(\mathcal{C}_{\mu}\) , so is \(A_{\lambda} \cap X_{\mu}\) . This completes the proof.

We remark that if the union of the \(\mathbf{X}_{\lambda}\) is different from \(\mathbf{X}\) , then the topology induced by \(\mathcal{G}\) on \(\mathbf{X} - \left(\bigcup_{\lambda \in \mathbf{L}} \mathbf{X}_{\lambda}\right)\) is discrete. For if \(x \in \mathbf{X}\) belongs to none of the \(\mathbf{X}_{\lambda}\) , then \(\{x\} \cap \mathbf{X}_{\lambda} = \emptyset\) is open in each topology \(\mathcal{G}_{\lambda}\) and therefore \(\{x\}\) is open in the topology \(\mathcal{G}\) .

